\section{Problemas durante as aulas}

Durante as aulas, foram propostos alguns problemas
introdutórios com o objetivo de revisar conceitos
fundamentais necessários ao desenvolvimento da
disciplina de Equações Diferenciais Ordinárias. Os
exercícios abordaram a resolução de equações
diferenciais simples a partir de integração indefinida,
bem como a determinação de soluções particulares
mediante condições iniciais. Esses problemas serviram
como uma revisão dos conceitos de derivação e
integração e como preparação para os exercícios da
primeira lista.

\question{
  Considere a equação diferencial:
  \[
    f'(x) + \sin{x} = 0.
  \]
  Determine \(y = f(x)\), sabendo que \(f(0) = 2\).
}

\answer{
  Partindo da equação diferencial,
  \begin{align*}
    f'(x) + \sin{x} &= 0 \\
    f'(x) &= -\sin{x}.
  \end{align*}

  Como \(f'(x) = \dfrac{dy}{dx}\), integra-se ambos os lados em relação a
  \(x\):
  \begin{align*}
    \frac{dy}{dx} &= -\sin{x} \\
    \int dy &= \int -\sin{x}\,dx \\
  \end{align*}

  Dessa forma conseguimos a solução geral:
  \begin{align*}
    y &= \cos{x} + C
  \end{align*}

  Aplicando a condição inicial \(f(0) = 2\),
  \begin{align*}
    f(0) &= \cos{0} + C \\
    2 &= 1 + C \\
    C &= 1
  \end{align*}

  Portanto, a solução particular é:
  \[
    \boxed{y = \cos{x} + 1}.
  \]
}

\newpage

\question{
  Encontre o valor de \(A\) para que
  \[
    y = A\sin{3t}
  \]
  seja solução da equação diferencial
  \[
    \frac{d^{2}y}{dt^{2}} + 2y = 4\sin{3t}
  \]
}

\answer{
  Calculando as derivadas de \(y = A\sin{3t}\), obtém-se
  \begin{align*}
    \frac{dy}{dt} &= 3A\cos{3t} \\
    \frac{d^{2}y}{dt^{2}} &= -9A\sin{3t}
  \end{align*}

  Substituindo \(y\) e sua segunda derivada na equação diferencial,
  \begin{align*}
    -9A\sin{3t} + 2A\sin{3t} &= 4\sin{3t} \\
    -7A\sin{3t} &= 4\sin{3t} \\
    -7A &= 4,
  \end{align*}

  portanto,
  \[
    \boxed{A = -\frac{4}{7}}.
  \]
}
